\chapter[Resultados: generación de reportes]{Resultados: generación de reportes y pipeline integrado}
\label{cp:res_gen}
% Extensión estimada: 25--30 páginas.
% Fuentes: sección "Generación de reportes" del documento de experimentación,
% Codigo/docs/resul/incertidumbre_llava_replica_resul_20260912.md y la sección "Costo computacional" del
% mismo documento.
%
% Feedback del tutor (14/09/2026): el capítulo de costo computacional era demasiado corto para sostenerse
% solo (3--5 páginas) y se incorpora como última sección de medición de este capítulo, que es donde se mide
% el pipeline completo. La comparación de hardware contra el estado del arte, que es interpretación y no
% medición, pasa al Capítulo~\ref{cp:discusion}. La sección de cierre se llama "Discusión".

\pendiente{Este capítulo responde la PI2 y evalúa el OE2, junto con la parte del OE3 que corresponde a la
generación de reportes (calidad del texto, tratamiento de la incertidumbre y costo computacional). Los reportes se generaron
sobre el test original, que es idéntico al del split por lesión; la robustez frente a la fuga de información
se reporta sobre el subconjunto limpio.}

\section{Generación directa con VLMs ajustados}
\label{sec:res_gen_directa}
\pendiente{Se presentan los resultados de Nivel 1 y Nivel 2 de los generadores ajustados sin
\emph{grounding}, que producen descripciones ricas pero diagnósticamente poco fieles.}

\section{Reporte frente a razonamiento}
\label{sec:res_gen_razonamiento}
\pendiente{Se describe el hallazgo metodológico sobre el entrenamiento con narrativas de razonamiento
(SkinCoT) y el \emph{trade-off} de estilo que distorsiona las métricas con referencia.}

\section{Ensamble en cascada}
\label{sec:res_gen_cascada}
\pendiente{Se justifica el uso de un ensamble y se presentan sus resultados con distintas fuentes de
diagnóstico (clasificador, DermLIP y \emph{oracle} como cota superior), con dos generadores y por clase.}

\section{Taxonomía de errores}
\label{sec:res_gen_errores}
\pendiente{Se clasifican los errores de los reportes (diagnóstico incorrecto, insuficiente, repetición, otros)
y se analiza cómo cambian con cada arquitectura.}

\section{Un generador dermatológico: SkinVL}
\label{sec:res_gen_skinvl}
\pendiente{Se evalúa si un generador especializado en dermatología evita la necesidad de la cascada.}

\section{Aporte de la imagen y del VLM}
\label{sec:res_gen_ablaciones}
\pendiente{Las ablaciones sin imagen o con imagen cambiada muestran que la imagen no cambia qué diagnóstico se
redacta, pero sí qué se describe. Se analiza qué aporta el VLM frente a una plantilla (riqueza clínica,
especificidad por caso, lenguaje): el clasificador aporta el diagnóstico y el VLM, la descripción.}

\section{Grounding por recuperación}
\label{sec:res_gen_rag}
\pendiente{Se compara el contexto construido con vecinos visuales frente al \emph{grounding} por
clasificador.}

\section{Ensamble consciente de la incertidumbre}
\label{sec:res_gen_incertidumbre}
\pendiente{Se presentan la distribución de casos por nivel (afirmativo, diferencial, abstención) y la
fidelidad diagnóstica sobre los casos no derivados y sobre la cohorte completa, con intervalos, junto con la
réplica con LLaVA-Med. Se explica cómo leer el valor de consistencia, que no es una exactitud multiclase
porque incluye el diferencial y excluye los casos derivados. Se reporta la robustez con el clasificador y el
kNN del split por lesión y sobre el subconjunto limpio.}

\section{Ejemplos cualitativos}
\label{sec:res_gen_cualitativos}
\pendiente{Se muestran reportes seleccionados que ilustran aciertos, diferenciales, derivaciones y errores
típicos.}

% Sección condicionada a que se concrete la evaluación con dermatólogos (septiembre--octubre 2026). Si no se
% realiza, se elimina y queda como limitación.
% \section{Evaluación por especialistas}
% \label{sec:res_gen_especialistas}

\section{Costo computacional del pipeline}
\label{sec:res_costo}
\pendiente{Se reportan el tiempo y la memoria de entrenamiento o ajuste de cada modelo en las GPUs
utilizadas, y el tiempo y la memoria por imagen de cada componente del \emph{pipeline} (clasificador, kNN y
generador) junto con el costo por reporte del pipeline completo. Estos números cierran el eje transversal de
hardware accesible planteado en la introducción; su comparación con los requerimientos de los trabajos del
estado del arte se retoma en el Capítulo~\ref{cp:discusion}.}

\section{Discusión}
\label{sec:res_gen_discusion}
\pendiente{Se resume qué arquitectura resulta más fiel diagnósticamente y bajo qué condiciones.}
